\section{Structure of higher-genus GW invariants of the quintic (Shuai Guo)}%
\label{sec:Structure of higher-genus GW invariants of the quintic (Shuai Guo)}

\subsection{The problem}%
\label{sub:The problem (Guo)}

We will consider the quintic threefold \(X = X_5 \subset \P^4\) and ask how many curves of genus \(g\) and degree \(d\) it contains. In genus zero, the answers begin
\[ 2875, \qquad 609250, \qquad 317206375, \qquad \ldots \]

Because \(X\) is Calabi-Yau, the moduli space \(\Mbar_{g,n}(X,\beta)\) has virtual dimension zero, and the string equation gives \(\ab<\1, \gamma_2, \ldots, \gamma_n>_{g,n,\beta} = 0\) unless \((g,n,\beta) = (0,3,0)\). We therefore define the primary invariants by
\[ N_{g,\beta} \coloneqq \ab<>_{g,0,\beta} = \int_{[\Mbar_g(X,\beta)]^{\vir}} 1 \]
and we package these numbers in the generating series
\[ F_g(Q) \coloneqq \sum_{\beta \text{ effective}} N_{g,\beta}Q^{\beta}. \]

The invariants \(N_{g,\beta}\) are rational numbers, while the integers we actually want are the BPS invariants \(n_{g,\beta}\). These are defined from the positive-degree part \(F_g^+\) by the multicover transformation
\[ \sum_{g \geq 0} F_g^+(Q)\lambda^{2g-2} = \sum_{\substack{g \geq 0 \\ \beta > 0}} n_{g,\beta} \sum_{k \geq 1} \frac{1}{k} \ab(2 \sin \frac{k\lambda}{2})^{2g-2}Q^{k\beta}. \]
Note that only the \(n_{g,\beta}\) are expected to be integral, and that the numbers above are the \(n_{0,d}\).

\begin{table}[htpb]
    \centering
    \caption{Leading BPS invariants \(n_{g,d}\) of the quintic threefold}
    \label{tab:quinticbps}
    \begin{tabular}{crrrrr}
        \toprule
        \(g\) & \(d=1\) & \(d=2\) & \(d=3\) & \(d=4\) & \(d=5\) \\
        \midrule
        0 & 2875 & 609250 & 317206375 & 242467530000 & 229305888887625 \\
        1 & 0 & 0 & 609250 & 3721431625 & 12129909700200 \\
        2 & 0 & 0 & 0 & 534750 & 75478987900 \\
        3 & 0 & 0 & 0 & 8625 & \(-15663750\) \\
        4 & 0 & 0 & 0 & 0 & 49250 \\
        5 & 0 & 0 & 0 & 0 & 1100 \\
        6 & 0 & 0 & 0 & 0 & 10 \\
        7 & 0 & 0 & 0 & 0 & 0 \\
        \bottomrule
    \end{tabular}
\end{table}

We see that every column of~\Cref{tab:quinticbps} eventually vanishes. This is the Castelnuovo bound, which for the quintic states that \(n_{g,d} = 0\) once \(g > \frac{d^2+5d+10}{10}\) and was proved by Liu--Ruan~\cite{LiuRuanCastelnuovo}. Integrality of the entries is the Gopakumar-Vafa conjecture, which was proved by Ionel--Parker~\cite{IonelParker2018}.

In genus zero, Candelas-de la Ossa-Green-Parkes computed \(F_0\) using period integrals on the mirror, which we discussed in~\Cref{sub:Enumerative mirror symmetry in genus zero}. In higher genus, Bershadsky-Cecotti-Ooguri-Vafa gave a recursion which produces \(F_g\) from the \(F_h\) with \(h < g\) up to finitely many unknown constants. Our goal is to explain how log GLSM determines this ambiguity.

\subsection{Genus zero and the polynomial structure of the B-model}%
\label{sub:Genus zero and the polynomial structure of the B-model}

Recall from~\Cref{sub:Enumerative mirror symmetry in genus zero} that for the quintic we consider the \(I\)-function
\[
    I(q,z) = z\sum_{d \geq 0} q^{H/z+d} \frac{\prod_{m=1}^{5d}(5H+mz)}{\prod_{m=1}^d (H+mz)^5} = I_0 z + I_1 H + I_2 \frac{H^2}{z} + I_3 \frac{H^3}{z^2},
\]
whose components are annihilated by the Picard-Fuchs operator
\[ D^4 - 5q(5D+1)(5D+2)(5D+3)(5D+4), \qquad D = q \odv{}{q}. \]
If we write \(J_k = I_k/I_0\) and \(\frac{I_1}{I_0} = \log q + \tau(q)\), so that \(Q = e^{I_1/I_0} = qe^{\tau(q)} = q + O(q^2)\), then the genus zero mirror theorem states that
\[ F_0(Q) = \frac{5}{2}\ab(J_1 J_2 - J_3). \]

For our purposes, it is useful to place the quintic formula in the general mirror-symmetry picture. Let \(\mc{X} = \ab\{(X,\omega)\}\) be a family of Calabi-Yau threefolds with Kähler moduli \(\mc{M} \ni Q = e^{2\pi i \omega}\) and let \(\mc{X}^{\vee} = \ab\{(X^{\vee},J)\}\) be a mirror family with complex moduli \(\mc{M}^{\vee} \ni q\). The solutions of the Picard-Fuchs equation are periods on \(X^{\vee}\) and assemble into
\[ I(z) = zI_0\varphi_0 + I_1 + I_2 z^{-1} + I_3 z^{-2}, \]
which lies in the span of \(z\varphi_0, \varphi_{\alpha}, \varphi^{\alpha}z^{-1}, \varphi^0 z^{-2}\). Here, \(I_1 = \sum_{\alpha} I_1^{\alpha}\varphi_{\alpha}\), \(I_2 = \sum_{\alpha} I_{2,\alpha}\varphi^{\alpha}\), and \(I_3 = I_{3,0}\varphi^0\), and the mirror map is again \(Q = e^{I_1/I_0} = qe^{\tau(q)}\).

We will use three special points on the complex moduli line of the quintic: the large complex structure point \(q=0\), the conifold point \(q = 5^{-5}\) (where a period vanishes), and the orbifold (or Landau-Ginzburg) point \(q = \infty\). The mirror map is defined at \(q=0\) and the expansion there produces Gromov-Witten invariants, while in higher genus we will impose boundary conditions at the other two points.

We now restate the mirror theorem in terms of three-point functions. Pick local coordinates \(q_i\) on \(\mc{M}^{\vee}\), let \(\ab\{\phi_i\}\) be the corresponding basis on the A-side, and define
\[ C^A_{ijk}(Q) \coloneqq I_0^2 \ab<\phi_i, \phi_j, \phi_k>^X_{0,3}, \qquad C^B_{ijk}(q) \coloneqq \int_{X^{\vee}} \partial_i \partial_j \partial_k \Omega \wedge \Omega. \]

\begin{thm}[Mirror theorem {\cite{Givental1998,LLY1997}}]
    We have \(C^A_{ijk}(Q(q)) = C^B_{ijk}(q)\). Moreover \(C^B_{ijk} \in \Delta^{-1}\Q[q]\), so it is a globally defined meromorphic function on \(\mc{M}^{\vee}\) with poles only along the discriminant.
\end{thm}

In particular, after multiplying by the correct power of \(I_0\) and applying the mirror map, a genus-zero generating function becomes a rational function of \(q\) with poles only along the discriminant. Therefore, it is determined by finitely many coefficients.

\begin{exm}
    For the quintic in the hyperplane basis,
    \[ C_{111} = \frac{5}{1-5^5 q}. \]
\end{exm}

\begin{exm}
    For \(X = X_{3,3} \subset \P^2 \times \P^2\), the Yukawa couplings are
    \begin{align*}
        C_{111} &= \frac{3^4 q_1 (2+3^3q_1+3^3q_2)}{\Delta}, & C_{222} &= \frac{3^4 q_2(2+3^3q_1+3^3q_2)}{\Delta}, \\
        C_{112} &= \frac{(1-3^3q_1-3^3q_2)(1+2\cdot 3^3 q_1 - 3^3 q_2)}{\Delta}, \\
        C_{122} &= \frac{(1-3^3q_1-3^3q_2)(1-3^3q_1+2\cdot3^3q_2)}{\Delta},
    \end{align*}
    where
    \[ \Delta = 1 - 3\cdot3^3(q_1+q_2) + 3\cdot3^6(q_1^2 - 7q_1q_2+q_2^2) - 3^9(q_1+q_2)^3. \]
\end{exm}

In higher genus, we expect normalized potentials \(f^A_{g,i}\) satisfying
\[ f^A_{g,i}(Q(q)) = f^B_{g,i}(q) \in \Delta^{-(2g-2+n)}\Q[q], \]
where the right hand side is computed by a sum over graphs.

\subsection{The BCOV Feynman rule}%
\label{sub:The BCOV Feynman rule}

The B-model Feynman rule and holomorphic anomaly equation originate in~\cite{BCOV1994}. For the quintic, Chang--Guo--Li prove polynomiality in~\cite{CGL2021} and the Feynman rule in~\cite{CGL2026}.

We will now discuss the BCOV recursion as a Feynman rule. It comes from a finite-dimensional model for the quantization of a symplectic transformation as in~\Cref{sub:Operator formalism and geometric quantization}. Let \(\mc{H} = \mc{H}^- \oplus \mc{H}^+ \subset \mc{H}_{\infty}\) be symplectic with \(d+1 = \dim \mc{H}^+ = \dim \mc{H}^-\) and set
\[ R = \begin{pmatrix} A & B \\ C & D \end{pmatrix} \in \End \mc{H}, \qquad \mc{Q}(\mbf{x},\bp) = (\bp, D^{-1}\mbf{x}) - \frac{1}{2}(\bp, D^{-1}C\bp). \]
Given a total potential \(F(\hbar, \mbf{x})\), we formally define the transformed potential by the Gaussian integral
\[ f(\hbar, \mbf{x}) \coloneqq (\wh{R}F)(\hbar,\mbf{x}) = \ln \int e^{\hbar^{-1}(\mc{Q}(\mbf{x},\bp') - \mbf{x}'\bp') + F(\hbar,\mbf{x}')}\d{\mbf{x}'}\d{\bp'}. \]
Expanding by stationary phase gives a sum over graphs. Here, \(D^{-1}C\) supplies the propagator on each edge and \(F\) supplies the vertices.

For a Calabi-Yau threefold the propagators are written \(S^{\varphi\varphi}\), \(S^{\varphi}\), and \(S\) according to the number of indices, and the vertices are the normalized A-model correlators
\[ P^A_{g,n}(Q) = I_0^{2-2g}\ab<\phi_1, \ldots, \phi_n>^X_{g,n}. \]
Note that these have nothing to do with the operator \(S_{\tau}(z)\) solving the quantum differential equation, which will appear later.

\begin{exm}
    For the quintic in genus two with no markings, the Feynman rule reads
    \begin{align*}
        f_2^B ={}& P_2 + \frac{1}{2}S^{\varphi\varphi}P_{1,1}^2 + \frac{1}{2}S^{\varphi\varphi}P_{1,2} + (S^{\varphi\varphi})^2 P_{1,1} + \frac{1}{8}(S^{\varphi\varphi})^2 P_{0,4} \\
        &+ \frac{1}{8}(S^{\varphi\varphi})^3 + \frac{1}{12}(S^{\varphi\varphi})^3 + \frac{\chi}{24}S^{\varphi}P_{1,1} + \frac{1}{2}\frac{\chi}{24}S^{\varphi}S^{\varphi\varphi} \\
        &+ \frac{1}{2}\frac{\chi}{24}\ab(\frac{\chi}{24}-1)S,
    \end{align*}
    where \(\chi = \chi(X_5) = -200\). This is a cubic polynomial in the propagators.
\end{exm}

We can interpret this as a sum over stable graphs, as in~\Cref{sub:R-matrix action}. A vertex \(v\) contributes the lower-genus correlator \(P_{g_v, \on{val}(v)}\); an internal edge contracts two of the insertions against a propagator; the legs record the remaining insertions. Each graph is weighted by \(\frac{1}{\ab|\Aut \Gamma|}\). As usual
\[ g(\Gamma) = h^1(\Gamma) + \sum_v g_v, \qquad 2g_v - 2 + \on{val}(v) > 0, \]
and \(f_g^B\) is the sum of the contributions of all stable graphs of genus \(g\).

\subsection{The holomorphic anomaly equation}%
\label{sub:The holomorphic anomaly equation}

We will now rewrite the recursion as a differential equation. The B-model amplitudes \(\mc{F}_g(z,\bar{z})\) are not holomorphic, and for \(g \geq 2\) their failure to be holomorphic is given by
\[ \partial_{\bar{i}} \mc{F}_g = \frac{1}{2}\bar{C}_{\bar{i}}^{jk}\ab( D_j D_k \mc{F}_{g-1} + \sum_{r=1}^{g-1} D_j \mc{F}_r D_k \mc{F}_{g-r} ). \]
Here, \(C_{ijk}\) is the Yukawa coupling, \(\bar{C}_{\bar{i}}^{jk}\) is obtained from \(\bar{C}_{\bar{i}\bar{j}\bar{k}}\) by raising two indices with the Weil-Petersson metric, and \(D_j\) is the covariant derivative on the appropriate power of the Hodge line bundle. The two terms on the right are the two boundary strata of \(\Mbar_g\), namely a non-separating and a separating node. The \(\frac{1}{2}\) forgets the ordering of the two branches.

If we choose local \(\bar{\partial}\)-primitives
\[ \partial_{\bar{i}}S^{jk} = \bar{C}_{\bar{i}}^{jk}, \qquad \partial_{\bar{i}}S^j = G_{i\bar{\ell}}S^{\ell j}, \qquad \partial_{\bar{i}}S = G_{i\bar{\ell}}S^{\ell}, \]
so that, modulo a holomorphic term, the anomaly equation becomes the polynomial PDE
\[ \pdv{\mc{F}_g}{S^{jk}} = \frac{1}{2}\ab( D_jD_k \mc{F}_{g-1} + \sum_{r=1}^{g-1} D_j \mc{F}_r D_k \mc{F}_{g-r} ). \]
These propagators are only defined up to holomorphic tensors, and this freedom is the source of the holomorphic ambiguity.

\begin{rmk}
    In the one-parameter normalized variables we will use for the quintic, this reads
    \[ -\partial_0 P_g = \frac{1}{2}P_{g-1,2} + \frac{1}{2}\sum_{r=1}^{g-1}P_{r,1}P_{g-r,1}, \qquad -\partial_1 P_g = 0. \]
\end{rmk}

To see why the graph sum satisfies this equation, note that in an unmarked potential the non-holomorphic generators occur only in the edge factors. A derivative in a propagator direction therefore picks out a single edge and cuts it. The two half-edges become markings on the resulting lower-genus vertices, and the covariant derivatives correspond to those markings by the divisor equation. We obtain \(P_{g-1,2}\) by cutting a loop and \(P_{g_1,1}P_{g_2,1}\) by cutting a separating edge.

\begin{rmk}
    For clarity, the implication only goes one way. The BCOV graph rule implies the anomaly equation, but the anomaly equation by itself does not determine the graph rule. To recover the same solution, we also need the propagator convention, the lower-genus vertex data, and the boundary conditions.
\end{rmk}

If all lower-genus vertices are known, then the right hand side is known, so we can integrate with respect to the non-holomorphic generators to obtain
\[ P_g = P_g^{\mr{particular}} + f_g(Y), \]
where finite generation gives the holomorphic ambiguity
\[ f_g(Y) = P_g|_{A = B_1 = B_2 = B_3 = 0} = \sum_{i=0}^{3g-3}a_{i,g}Y^i, \qquad Y = (1-5^5q)^{-1}. \]
We see that the anomaly equation leaves \(3g-2\) unknown coefficients in each genus.

\begin{rmk}
    For the quintic, the degree-zero data together with orbifold regularity remove \(\ceil{\frac{3g-3}{5}}+1\) of these freedoms and the conifold gap supplies \(2g-2\) further conditions, so that
    \[ 3g-2 - \ab(\ceil*{\frac{3g-3}{5}}+1) - (2g-2) = \floor*{\frac{2g-2}{5}} \]
    parameters remain and are genuinely new effective invariants. This is exactly the number of \(\infty\)-vertex unknowns predicted by~\Cref{eqn:balancing condition}, namely the number of \(\beta\) with \(\beta \leq \frac{2g-2}{5}\).
\end{rmk}

We therefore obtain the following conjectural algorithm proposed by physicists: the BCOV Feynman rule, orbifold regularity at the Landau-Ginzburg point (a form of the LG/CY correspondence, which for the quintic is work in progress of Guo-Ruan-Janda), and the conifold gap condition. To these we may add the Gopakumar-Vafa integrality and the Castelnuovo bound.

\subsection{Four theories}%
\label{sub:Four theories}

Recall from~\Cref{sec:Log GLSM (Qile Chen)} that the log GLSM carries two virtual cycles. The reduced cycle recovers quintic Gromov-Witten theory, while the canonical equivariant cycle admits a root of unity specialization which is semisimple and therefore computable by Givental-Teleman. The bipartite graph sum
\[ \Omega^{\mr{GW}}_{g,n} = \sum_{\Gamma \text{ bipartite}} \frac{1}{\ab|\Aut \Gamma|} \ab( \prod_{v \in V_{\infty}(\Gamma)} c^{\mr{eff}}_{g_v,\beta_v}) \Cont^{\mr{can}}_{\Gamma} \]
expresses the former in terms of the latter, at the cost of introducing the effective invariants \(c^{\mr{eff}}_{g,\beta}\) at the \(\infty\)-vertices. Replacing these by formal parameters \(\mbf{C} = \ab\{c_{g,\beta}\}\) gives the \textit{extended \(\mbf{C}\)-family}, whose vertex data is supplied by the meromorphic and rubber log GLSM theories.

We will consider the following four theories.
\begin{itemize}
    \item The \textit{\(t\)-twisted geometric theory} sets \(\lambda_i = 0\) and keeps the auxiliary weight \(t\). It has both stable-map and stable-quotient versions, and the quotient version gives the genus two localization formula below.
    \item The \textit{formal quintic}, or formal \(\lambda\)-twisted theory, sets \(t = 0\) and \(\lambda_i = \zeta_5^i \lambda\). It is a semisimple CohFT, so it has a universal and computable \(R\)-matrix.
    \item The \textit{extended \(\mbf{C}\)-family} interpolates between these, and is where bipartite graphs with both formal and meromorphic vertices live.
    \item \textit{Quintic Gromov-Witten theory} is obtained by setting \(\mbf{C} = \mbf{C}^{\mr{eff}}\) and then letting \(\lambda \to 0\), and is the theory we want to compute.
\end{itemize}

\begin{warn}
    In higher genus the \(t^0\) coefficient of the \(t\)-twisted stable quotient theory is not the \(\lambda^0\) coefficient of the formal \(\lambda\)-twisted theory.
\end{warn}

\subsection{The formal quintic}%
\label{sub:The formal quintic}

The formal quintic and its holomorphic anomaly equations are studied in~\cite{LhoPandharipande2019}; the comparison with the extended theory is developed in~\cite{GJRstructure}.

Let \(E = \mc{O}_{\P^4}(5)\) and let \(\pi \colon \mc{C} \to \Mbar_{g,n}(\P^4,d)\) be the universal curve. The \(t\)-twisted cycle is
\[ e(R\pi_* f^* E \otimes [1]) \cap [\Mbar_{g,n}(\P^4,d)]^{\vir}, \]
where \([1]\) has equivariant first Chern class \(t\). Starting instead from the fully \((\C^{\times})^5\)-equivariant cycle, the formal quintic is the root of unity specialization
\[ [\Mbar_{g,n}(\P^4,d)]^{\vir,\lambda} \coloneqq \ab( e(R\pi_*f^*E \otimes [1]) \cap [\Mbar_{g,n}(\P^4,d)]^{\vir}_{(\C^{\times})^5})\Big|_{t=0, \; \lambda_i = \zeta_5^i \lambda}, \]
where \(\zeta_5 = e^{2\pi i/5}\). The relation \(H^5 = \lambda^5\) makes the quantum product generically semisimple while retaining \(\lambda\) as a parameter.

The state space, unit, and pairing are given by
\[ \mc{H}_{\lambda} = \Q(\lambda)[H]/(H^5-\lambda^5), \quad \1 = 1, \quad (\gamma_1,\gamma_2)^{\lambda} = \int_{\P^4, (\C^{\times})^5} e(\mc{O}(5))\gamma_1\gamma_2 \Big|_{\lambda_i = \zeta_5^i\lambda}, \]
and the CohFT classes are
\[ \Omega^{\lambda}_{g,n}(\gamma_1, \ldots, \gamma_n) = \sum_{d \geq 0} q^d \rho_* \ab( \prod_{i=1}^n \ev_i^*(\gamma_i) \cap [\Mbar_{g,n}(\P^4,d)]^{\vir,\lambda}), \]
where \(\rho\) forgets the map and stabilizes the domain. If we shift by \(\tau(q)H\) with \(\tau(q) = I^{\lambda}_1(q)/I^{\lambda}_0(q)\), then the shifted theory \(\Omega^{\lambda,\tau}\) is semisimple. After extending scalars to \(\Q(\zeta_5,\lambda)\), the idempotent basis satisfies
\[ \omega^{\lambda}_{g,n}(e_{\alpha_1}, \ldots, e_{\alpha_n}) = \sum_{\alpha}\Delta_{\alpha}^{g-1}\prod_i \delta_{\alpha,\alpha_i}. \]
Then~\Cref{thm:teleman} gives
\[ \Omega^{\lambda,\tau}_{g,n} = \sum_{\Gamma \in G_{g,n}} \frac{1}{\ab|\Aut\Gamma|}\iota_* \Cont_{\Gamma} \]
with the decorations of~\Cref{sub:R-matrix action}. A leg carrying \(\gamma\) gets \(R^{-1}(\bar{\psi})\gamma\), an edge gets
\[ V_R(z_1,z_2) = \frac{\sum_{\alpha} e_{\alpha}\otimes e^{\alpha} - \sum_{\alpha} R^{-1}(z_1)e_{\alpha} \otimes R^{-1}(z_2)e^{\alpha}}{z_1+z_2}, \]
and a vertex gets the topological field theory \(\omega\), with the missing unit-preserving part inserted at forgotten markings via \(T(z) = z(\1 - R^{-1}(z)\1)\).

\begin{rmk}
    We see that the edge factor has no pole. Setting \(z_2 = -z_1\) and using the symplectic condition in the form \(R^{-1}(z) = R^*(-z)\), the numerator vanishes. Therefore, it is divisible by \(z_1+z_2\), and \(V_R\) is a power series in the two cotangent line variables. The coefficient of \(z_1^p z_2^q\) is the class placed at a node with \(\bar{\psi}_1^p \bar{\psi}_2^q\), and contracting it against the two adjacent topological field theory tensors glues the branches.
\end{rmk}

\subsection{How to compute the \(R\)-matrix}%
\label{sub:How to compute the R-matrix (Guo)}

To compute the \(R\)-matrix, we begin with the formal equivariant \(I\)-function
\[ I^{\lambda}(q,z) = z\sum_{d\geq0}q^d \frac{\prod_{m=1}^{5d}(5H+mz)}{\prod_{m=1}^d ((H+mz)^5 - \lambda^5)} = zI_0 + I_1^{\lambda}H + z^{-1}I_2^{\lambda}H^2 + \cdots. \]
The coefficients through \(H\) are independent of \(\lambda\), so \(I_1^{\lambda} = I_1 - I_0 \log q\) and the mirror shift \(\tau(q) = I_1^{\lambda}(q)/I_0(q)\) is free of logarithms. Set
\[ D = q\odv{}{q}, \qquad Y = (1-5^5q)^{-1}, \qquad L = Y^{1/5} \]
and introduce the bases
\[ \phi_k = I_{1,1}\cdots I_{k,k}H^k, \qquad \bar{\phi}_k = I_0 L^{-k-1}\phi_k, \]
where the diagonal Birkhoff factors are defined recursively by
\[ I_{p,0} = I_p, \qquad I_{p,r} = \frac{I_{p-1,r-1}}{I_{r-1,r-1}} + D \ab( \frac{I_{p,r-1}}{I_{r-1,r-1}}), \quad 1 \leq r \leq p, \]
so that in particular \(I_{1,1} = 1+D\tau\). The idempotents and canonical coordinates are given by
\[ e_{\alpha} = \frac{1}{5}\sum_{i=0}^4 h_{\alpha}^{-i}\frac{\phi_i}{L^i}, \qquad u_{\alpha} = h_{\alpha}\int_0^q (L(s)-1)\frac{\d{s}}{s}, \qquad h_{\alpha} = \zeta_5^{\alpha}\lambda. \]

We will use the quantum differential equation to determine the \(q\)-dependence of the \(R\)-matrix. The fundamental solution satisfies
\[ zDS_{\tau}(z) + S_{\tau}(z)H = M_{\tau}S_{\tau}(z), \qquad M_{\tau} = \dot{\tau}\star, \quad \dot{\tau} \coloneqq (1+D\tau)H = I_{1,1}H, \]
and in the normalized canonical direction \(\alpha\) the stationary phase asymptotics give
\[ (S_{\tau})_{i\bar{\alpha}}(z)C_{\alpha}(z)^{-1} = e^{u_{\alpha}/z}R_{i\bar{\alpha}}(z), \]
so that
\[ (zD+h_{\alpha}L)R_{i\bar{\alpha}} = \sum_j (M_{\tau})_i^j R_{j\bar{\alpha}}. \]
Here, \(C_{\alpha}(z)\) is the \(q=0\) quantum Riemann-Roch (equivalently, Gamma function) factor, and together with the symplectic condition it fixes the constants that the differential equation cannot see. As at the end of~\Cref{sub:How to compute the R-matrix}, the differential equation gives a recursion and an oscillating integral gives the initial conditions.

Writing \(R(z) = \mr{Id} + R_1 z + R_2 z^2 + \cdots\), we can use the equation above to determine each \(R_{k+1}\) from the preceding ones. There are two features which make the quintic tractable. First, the root of unity basis has five sectors, and a mod \(5\) charge rule permits only one residue pattern in each row. Then most matrix entries vanish before any computation starts. Second, every coefficient lies in the finitely generated ring
\[ \mc{R} = \Q[A,B_1,B_2,B_3,Y], \qquad A = \frac{DI_{1,1}}{I_{1,1}}, \quad B_k = \frac{D^k I_0}{I_0}, \]
so the stable graph sum reduces to a finite algebraic computation.

\begin{rmk}
    In particular, a fixed genus only needs finitely much of \(R\) for dimensional reasons. On \(\Mbar_{g,n}\) only classes of total degree at most \(3g-3+n\) contribute and everything else integrates to zero, so for an unmarked potential we need only \(R_0, \ldots, R_{3g-3}\). Genus \(4\) stops at \(z^9\) and genus \(5\) at \(z^{12}\).
\end{rmk}

